\section{Introduction}
\label{sec:introduction}

Whether formulating a national policy or simply deciding what groceries to buy, we increasingly rely on the
charts, figures and text created by scientists and journalists. Interpreting these visual and textual
summaries is essential to making informed decisions. However, most of the artefacts we encounter are \emph{opaque}: they are unable to reveal anything about how they relate to the data they were derived from. While one could in principle try to use the source code and data sources to reverse engineer some of these relationships, this requires substantial expertise, as well as valuable time spent away from the ``comprehension context'' in which we encountered the output in question.  These difficulties are only compounded when the information presented draws on multiple data sources, such as medical meta-analyses~\citep{dersimonian86}, ensemble models~\citep{murphy04}, or queries that span multiple database tables~\citep{selinger79}. Even professional reviewers may lack the resources or inclination to get too involved. Perhaps more often than we would like, we end up taking things on trust.
